% ============================================================================
%  ESTADÍSTICA INFERENCIAL — Sesión 8
%  Semana 8: 28 de septiembre al 3 de octubre de 2026
%  2.3 Distribución de probabilidad continua: normal
% ============================================================================
\documentclass[aspectratio=169,10pt]{beamer}

\usepackage[utf8]{inputenc}
\IfFileExists{spanish.ldf}{\usepackage[spanish,es-noshorthands]{babel}}{}
\usepackage{amsmath,amssymb}
\usepackage{booktabs}
\usepackage{tikz}
\usepackage{listings}

\usetheme{Uniminuto}

\uniasignatura{Estadística inferencial}
\uniprofesor{Andrés Felipe Puerta Velez}
\unicorreo{andres.puerta-v@uniminuto.edu.co}
\uninrc{CBASBL061}
\unisemana{Sesión 8 — Semana 8 · 28 de septiembre al 3 de octubre de 2026}

\newcommand{\listatiny}{\setbeamerfont{itemize/enumerate body}{size=\tiny}\tiny}

% curva normal reutilizable
\newcommand{\curvanormal}[1]{%
  \draw[uniGris,->] (-3.4,0) -- (3.5,0);
  \draw[uniAzul,thick,domain=-3.3:3.3,samples=90,smooth]
    plot (\x,{2.1*exp(-0.5*\x*\x)});
  #1}

\newcommand{\fx}[1]{{\ttfamily\color{uniAzul}#1}}

% ------------------------------------------------------------- CÓDIGO R ----
\lstdefinestyle{Rcode}{
  language=R,
  basicstyle=\ttfamily\scriptsize,
  keywordstyle=\color{uniAzul}\bfseries,
  commentstyle=\color{uniGris}\itshape,
  stringstyle=\color{uniTitulo},
  showstringspaces=false, numbers=none, breaklines=true,
  columns=fullflexible,
  backgroundcolor=\color{uniCremaClaro},
  frame=leftline, framerule=1.2pt, rulecolor=\color{uniNaranja},
  xleftmargin=0.35cm, framexleftmargin=0.35cm,
  framexrightmargin=0.15cm, framextopmargin=2pt, framexbottommargin=2pt,
  aboveskip=0.15cm, belowskip=0.15cm
}
\lstdefinestyle{Rmini}{
  language=R,
  basicstyle=\ttfamily\tiny,
  keywordstyle=\color{uniAzul}\bfseries,
  commentstyle=\color{uniGris}\itshape,
  stringstyle=\color{uniTitulo},
  showstringspaces=false, numbers=none, breaklines=true,
  columns=fullflexible,
  backgroundcolor=\color{uniCremaClaro},
  frame=leftline, framerule=1.2pt, rulecolor=\color{uniNaranja},
  xleftmargin=0.35cm, framexleftmargin=0.35cm,
  framexrightmargin=0.15cm, framextopmargin=2pt, framexbottommargin=2pt,
  aboveskip=0.12cm, belowskip=0.12cm
}
\lstdefinestyle{Rout}{
  basicstyle=\ttfamily\tiny\color{black},
  showstringspaces=false, numbers=none, breaklines=true,
  columns=fullflexible,
  backgroundcolor=\color{uniGris!10},
  frame=leftline, framerule=1.2pt, rulecolor=\color{uniGris},
  xleftmargin=0.35cm, framexleftmargin=0.35cm,
  framexrightmargin=0.15cm, framextopmargin=2pt, framexbottommargin=2pt,
  aboveskip=0.12cm, belowskip=0.12cm
}

\begin{document}

\portada

\begin{frame}{Orden del día}
  \begin{enumerate}
    \item Llamado a lista y revisión de los ejercicios propuestos
    \item Repaso: dónde quedamos
    \item \azul{Tema 2.3: distribución normal}
    \item Estandarización y uso de la tabla $Z$
    \item Ejercicios en clase
    \item Cierre y ejercicios propuestos
  \end{enumerate}
\end{frame}

\begin{frame}{Repaso: dónde quedamos}
  En la sesión 7 modelamos variables \azul{discretas}: binomial (éxitos en $n$
  ensayos) y Poisson (ocurrencias en un intervalo). En ambas se pregunta
  $P(X=k)$ para valores aislados.

  \vspace{0.25cm}
  \begin{alertblock}{El cambio de hoy}
    \footnotesize
    Muchas variables no son contables sino \destacar{medibles}: tiempos, pesos,
    ingresos, longitudes. Pueden tomar cualquier valor en un intervalo. Para
    ellas $P(X=k)=0$ y solo tienen sentido preguntas del tipo
    $P(a<X<b)$: \azul{la probabilidad es área bajo la curva}.
  \end{alertblock}
\end{frame}

\seccion{2.3 Distribución normal}

\begin{frame}{La distribución normal}
  \begin{columns}[T,onlytextwidth]
    \begin{column}{0.46\textwidth}
      \footnotesize
      \begin{block}{Características}
        \begin{itemize}
          \item Forma de campana, simétrica respecto a $\mu$.
          \item Media, mediana y moda coinciden.
          \item Queda definida por dos parámetros: $\mu$ y $\sigma$.
          \item El área total bajo la curva es 1.
          \item Se extiende de $-\infty$ a $+\infty$.
        \end{itemize}
      \end{block}
      Se escribe $X\sim N(\mu,\sigma)$.
    \end{column}
    \begin{column}{0.52\textwidth}
      \centering
      \begin{tikzpicture}[scale=0.68]
        \curvanormal{}
        \draw[uniNaranja,dashed,thick] (0,0) -- (0,2.1);
        \node[font=\scriptsize] at (0,-0.3) {$\mu$};
        \node[font=\scriptsize] at (1,-0.3) {$+\sigma$};
        \node[font=\scriptsize] at (-1,-0.3) {$-\sigma$};
        \node[font=\scriptsize] at (2,-0.3) {$+2\sigma$};
        \node[font=\scriptsize] at (-2,-0.3) {$-2\sigma$};
      \end{tikzpicture}
    \end{column}
  \end{columns}

  \vspace{0.02cm}
  {\scriptsize\color{uniGris} Aparece por todas partes porque muchas
  variables son la suma de muchos efectos pequeños e independientes. La razón
  formal es el teorema del límite central, que veremos en la sesión 9.}
\end{frame}

\begin{frame}{Regla empírica: 68 -- 95 -- 99{,}7}
  \begin{columns}[T,onlytextwidth]
    \begin{column}{0.52\textwidth}
      \centering
      \begin{tikzpicture}[scale=0.68]
        \fill[uniAzul,opacity=0.30,domain=-1:1,samples=60]
          plot (\x,{2.1*exp(-0.5*\x*\x)}) -- (1,0) -- (-1,0) -- cycle;
        \fill[uniNaranja,opacity=0.22,domain=-2:-1,samples=40]
          plot (\x,{2.1*exp(-0.5*\x*\x)}) -- (-1,0) -- (-2,0) -- cycle;
        \fill[uniNaranja,opacity=0.22,domain=1:2,samples=40]
          plot (\x,{2.1*exp(-0.5*\x*\x)}) -- (2,0) -- (1,0) -- cycle;
        \curvanormal{}
        \node[font=\scriptsize] at (0,-0.3) {$\mu$};
        \node[font=\scriptsize] at (1,-0.3) {$1\sigma$};
        \node[font=\scriptsize] at (2,-0.3) {$2\sigma$};
        \node[font=\scriptsize] at (-1,-0.3) {$-1\sigma$};
        \node[font=\scriptsize] at (-2,-0.3) {$-2\sigma$};
      \end{tikzpicture}
    \end{column}
    \begin{column}{0.44\textwidth}
      \begin{block}{Áreas centrales}
        \footnotesize
        \begin{center}
        \begin{tabular}{@{}lr@{}}
          \toprule
          Intervalo & Área \\
          \midrule
          $\mu\pm 1\sigma$ & 68{,}27\% \\
          $\mu\pm 2\sigma$ & 95{,}45\% \\
          $\mu\pm 3\sigma$ & 99{,}73\% \\
          \bottomrule
        \end{tabular}
        \end{center}
      \end{block}
      {\footnotesize Un valor a más de $3\sigma$ de la media es \azul{muy}
      raro: menos de 3 en 1\,000.}
    \end{column}
  \end{columns}
\end{frame}

\begin{frame}{Estandarización: la variable $Z$}
  \begin{block}{Transformación}
    \[ Z=\frac{X-\mu}{\sigma} \qquad\Longrightarrow\qquad Z\sim N(0,1) \]
    $Z$ indica \destacar{a cuántas desviaciones estándar} está el valor $X$ de
    la media.
  \end{block}

  \vspace{0.2cm}
  \begin{columns}[T,onlytextwidth]
    \begin{column}{0.48\textwidth}
      \begin{exampleblock}{Por qué sirve}
        \footnotesize
        Hay infinitas normales, una por cada par $(\mu,\sigma)$. Pero todas se
        reducen a \azul{una sola} tabla: la de $Z$.
      \end{exampleblock}
    \end{column}
    \begin{column}{0.48\textwidth}
      \begin{alertblock}{Cómo se lee la tabla}
        \footnotesize
        La tabla da $P(Z<z)$, el área \azul{a la izquierda}. Por tanto:
        \[ P(Z>z)=1-P(Z<z) \]
        \[ P(a<Z<b)=P(Z<b)-P(Z<a) \]
      \end{alertblock}
    \end{column}
  \end{columns}
\end{frame}

\begin{frame}{Valores de $Z$ que conviene memorizar}
  \footnotesize
  \begin{center}
  \begin{tabular}{@{}crll@{}}
    \toprule
    $z$ & $P(Z<z)$ & \textbf{Área central} & \textbf{Uso} \\
    \midrule
    $1{,}00$  & $0{,}8413$ & 68{,}27\% & regla empírica \\
    $1{,}645$ & $0{,}9500$ & 90\% & confianza del 90\% \\
    $1{,}96$  & $0{,}9750$ & 95\% & confianza del 95\% \\
    $2{,}00$  & $0{,}9772$ & 95{,}45\% & regla empírica \\
    $2{,}576$ & $0{,}9950$ & 99\% & confianza del 99\% \\
    \bottomrule
  \end{tabular}
  \end{center}

  \vspace{0.25cm}
  \begin{alertblock}{Por simetría}
    \footnotesize
    $P(Z<-z)=1-P(Z<z)$. Por ejemplo $P(Z<-1)=1-0{,}8413=0{,}1587$. Las tablas
    que solo traen valores positivos se usan así.
  \end{alertblock}

  \vspace{0.1cm}
  {\scriptsize\color{uniGris} Estos son los mismos $z$ del tamaño de muestra de
  la sesión 5. Ahora se entiende de dónde salían.}
\end{frame}

\begin{frame}{Ejercicio 1 — en clase}
  \begin{alertblock}{Individual (9 minutos)}
    El tiempo de entrega de una empresa de domicilios se distribuye normal con
    media \destacar{$\mu=500$ minutos} y desviación \destacar{$\sigma=50$
    minutos}.

    \vspace{0.25cm}
    \begin{enumerate}
      \item $P(X<560)$
      \item $P(X<450)$
      \item $P(450<X<560)$
      \item $P(X>600)$
    \end{enumerate}
  \end{alertblock}
\end{frame}

\begin{frame}{Ejercicio 1 — solución}
  \small
  \begin{block}{Primero se estandariza}
    \begin{align*}
      x=560 &\Rightarrow z=\tfrac{560-500}{50}=1{,}20
        &&P(X<560)=\azul{0{,}8849}\\[0.2em]
      x=450 &\Rightarrow z=\tfrac{450-500}{50}=-1{,}00
        &&P(X<450)=\azul{0{,}1587}\\[0.2em]
      x=600 &\Rightarrow z=\tfrac{600-500}{50}=2{,}00
        &&P(X>600)=1-0{,}9772=\azul{0{,}0228}
    \end{align*}
  \end{block}

  \vspace{0.1cm}
  \begin{exampleblock}{Punto 3}
    \footnotesize
    $P(450<X<560)=0{,}8849-0{,}1587=\azul{0{,}7263}$: cerca del 73\% de los
    pedidos se entrega entre 450 y 560 minutos.
  \end{exampleblock}
\end{frame}

\begin{frame}{Ejercicio 2 — en clase}
  \begin{alertblock}{En parejas (9 minutos)}
    El peso neto de un producto empacado sigue una normal con
    \destacar{$\mu=2{,}5$ kg} y \destacar{$\sigma=0{,}3$ kg}. La norma exige
    que el contenido no baje de 2{,}0 kg.

    \vspace{0.25cm}
    \begin{enumerate}
      \item ¿Qué porcentaje de empaques incumple la norma?
      \item $P(2{,}2<X<2{,}8)$
      \item $P(X>3{,}1)$
      \item ¿Por debajo de qué peso está el 90\% de los empaques?
    \end{enumerate}
  \end{alertblock}
\end{frame}

\begin{frame}{Ejercicio 2 — solución}
  \begin{columns}[T,onlytextwidth]
    \begin{column}{0.52\textwidth}
      \begin{block}{Cálculos directos}
        \footnotesize
        \begin{align*}
          z&=\tfrac{2{,}0-2{,}5}{0{,}3}=-1{,}67 &&P(X<2{,}0)=\azul{0{,}0478}\\[0.2em]
          &\text{$2{,}2$ y $2{,}8$ están a $\pm 1\sigma$} &&P=\azul{0{,}6827}\\[0.2em]
          z&=\tfrac{3{,}1-2{,}5}{0{,}3}=2{,}00 &&P(X>3{,}1)=\azul{0{,}0228}
        \end{align*}
      \end{block}
    \end{column}
    \begin{column}{0.44\textwidth}
      \begin{exampleblock}{Punto 4: problema inverso}
        \footnotesize
        Se busca $x$ tal que $P(X<x)=0{,}90$. En la tabla, $z=1{,}28$:
        \[ x=\mu+z\sigma=2{,}5+1{,}28(0{,}3)=\azul{2{,}88\text{ kg}} \]
      \end{exampleblock}
      \begin{alertblock}{Ojo}
        \footnotesize
        Aquí se entra a la tabla \azul{por el área} y se sale con el $z$: es el
        camino contrario al de los puntos 1 a 3.
      \end{alertblock}
    \end{column}
  \end{columns}
\end{frame}

\begin{frame}{Los dos tipos de problema}
  \footnotesize
  \begin{center}
  \begin{tabular}{@{}p{4.6cm}p{3.2cm}p{3.6cm}@{}}
    \toprule
    \textbf{Pregunta} & \textbf{Camino} & \textbf{Ejemplo} \\
    \midrule
    Dado un valor, ¿qué probabilidad? &
      $x \rightarrow z \rightarrow$ tabla $\rightarrow$ área &
      ``¿qué \% pesa menos de 2 kg?'' \\
    \addlinespace
    Dada una probabilidad, ¿qué valor? &
      área $\rightarrow$ tabla $\rightarrow z \rightarrow x=\mu+z\sigma$ &
      ``¿bajo qué peso está el 90\%?'' \\
    \bottomrule
  \end{tabular}
  \end{center}

  \vspace{0.25cm}
  \begin{alertblock}{Recomendación práctica}
    \footnotesize
    \destacar{Dibuje siempre la campana} y sombree lo que le piden antes de
    buscar en la tabla. La mitad de los errores en el examen vienen de sombrear
    la cola equivocada.
  \end{alertblock}
\end{frame}

\begin{frame}[fragile]{La normal, en R: se acaba la tabla}
  \scriptsize
\begin{lstlisting}[style=Rmini]
# Ejercicio 1: X ~ N(500, 50)
pnorm(560, mean = 500, sd = 50)          # P(X < 560)
pnorm(450, mean = 500, sd = 50)          # P(X < 450)
1 - pnorm(600, mean = 500, sd = 50)      # P(X > 600)
diff(pnorm(c(450, 560), 500, 50))        # P(450 < X < 560)

# Ejercicio 2: X ~ N(2.5, 0.3)
pnorm(2.0, 2.5, 0.3)                     # incumple la norma
qnorm(0.90, 2.5, 0.3)                    # problema inverso: el percentil 90

# Los z que memorizamos
qnorm(c(0.95, 0.975, 0.995))
\end{lstlisting}
\begin{lstlisting}[style=Rout]
[1] 0.8849303
[1] 0.1586553
[1] 0.02275013
[1] 0.7262751
[1] 0.04779035
[1] 2.884466
[1] 1.644854 1.959964 2.575829
\end{lstlisting}

  \vspace{-0.12cm}
  \begin{alertblock}{Lo que cambia y lo que no}
    \tiny
    Con \fx{pnorm()} \destacar{ya no hace falta estandarizar}: se le pasan
    $\mu$ y $\sigma$ y devuelve el área. Pero el examen se resuelve con tabla,
    y sobre todo: \destacar{decidir si el problema es directo o inverso}
    —área o valor— sigue siendo trabajo suyo.
  \end{alertblock}
\end{frame}

\begin{frame}{Cierre de la sesión}
  \begin{block}{Lo que debe quedar claro hoy}
    \begin{itemize}
      \item En variables continuas la probabilidad es \azul{área}, y
            $P(X=k)=0$.
      \item $Z=(X-\mu)/\sigma$ convierte cualquier normal en la estándar.
      \item La tabla da el área a la izquierda; lo demás sale por complemento,
            resta o simetría.
      \item Problema inverso: de la probabilidad al valor, con
            $x=\mu+z\sigma$.
    \end{itemize}
  \end{block}

  \vspace{0.2cm}
  \begin{alertblock}{Trabajo independiente para la próxima sesión}
    \footnotesize
    Resolver los \destacar{20 ejercicios de la lámina siguiente}. Traer impresa
    la tabla de la normal estándar: la usaremos el resto del curso.
  \end{alertblock}
\end{frame}

\begin{frame}{Ejercicios propuestos}
  \vspace{-0.1cm}
  \begin{beamercolorbox}[wd=\textwidth,sep=0.10cm]{block body}
    \centering\scriptsize
    \textbf{1--8:} $Z\sim N(0,1)$. \quad \textbf{9--20:} $X\sim N(70,10)$.
  \end{beamercolorbox}
  \listatiny
  \vspace{0.15cm}
  \begin{columns}[T,onlytextwidth]
    \begin{column}{0.47\textwidth}
      \begin{enumerate}\setlength{\itemsep}{0.16cm}
        \item $P(Z<1{,}00)$
        \item $P(Z<1{,}50)$
        \item $P(Z<-1{,}00)$
        \item $P(Z>1{,}96)$
        \item $P(-1<Z<1)$
        \item $P(-2<Z<2)$
        \item $z$ tal que $P(Z<z)=0{,}95$
        \item $z$ tal que $P(Z<z)=0{,}99$
        \item $P(X<85)$
        \item $P(X<60)$
      \end{enumerate}
    \end{column}
    \begin{column}{0.50\textwidth}
      \begin{enumerate}\setlength{\itemsep}{0.16cm}
        \setcounter{enumi}{10}
        \item $P(X<50)$
        \item $P(X>90)$
        \item $P(60<X<85)$
        \item $P(50<X<90)$
        \item $z$ correspondiente a $x=85$
        \item $z$ correspondiente a $x=50$
        \item Percentil 95 de $X$
        \item Percentil 25 de $X$
        \item Intervalo $\mu\pm 2\sigma$
        \item \% de datos dentro de $\mu\pm 3\sigma$
      \end{enumerate}
    \end{column}
  \end{columns}
\end{frame}

\begin{frame}{Respuestas}{Autoevaluación de los ejercicios propuestos}
  \listatiny
  \vspace{0.15cm}
  \begin{columns}[T,onlytextwidth]
    \begin{column}{0.47\textwidth}
      \begin{enumerate}\setlength{\itemsep}{0.16cm}
        \item $0{,}8413$
        \item $0{,}9332$
        \item $0{,}1587$
        \item $0{,}0250$
        \item $0{,}6827$
        \item $0{,}9545$
        \item $1{,}645$
        \item $2{,}326$
        \item $0{,}9332$
        \item $0{,}1587$
      \end{enumerate}
    \end{column}
    \begin{column}{0.50\textwidth}
      \begin{enumerate}\setlength{\itemsep}{0.16cm}
        \setcounter{enumi}{10}
        \item $0{,}0228$
        \item $0{,}0228$
        \item $0{,}7745$
        \item $0{,}9545$
        \item $z=1{,}5$
        \item $z=-2{,}0$
        \item $86{,}45$
        \item $63{,}26$
        \item $[50;\,90]$
        \item $99{,}73\%$
      \end{enumerate}
    \end{column}
  \end{columns}
\end{frame}

\begin{frame}{Referencias}
  \footnotesize
  \begin{itemize}
    \item[{[1]}] Díaz Rodríguez, M. (2022). \textit{Estadística inferencial
      aplicada}. Universidad del Norte.
      \texttt{https://elibro.net/es/lc/uniminuto/titulos/221614}
    \item[{[2]}] Llinás Solano, H. (2014). \textit{Introducción a la
      estadística con aplicaciones en ciencias sociales}. Universidad del
      Norte. \texttt{https://elibro.net/es/lc/uniminuto/titulos/70062}
    \item[{[3]}] Contento Rubio, M. R. (2019). \textit{Estadística con
      aplicaciones en R}. Editorial Utadeo.
      \texttt{https://elibro.net/es/lc/uniminuto/titulos/220926}
  \end{itemize}

  \vspace{0.35cm}
  \begin{center}
    {\Large\color{uniAzul}\textbf{¿Preguntas?}}
  \end{center}
\end{frame}

\end{document}
